\documentclass[tikz, border=6pt]{standalone}
\usepackage[fontset=fandol]{ctex}
\usetikzlibrary{positioning, arrows.meta, calc, fit, backgrounds}

\begin{document}
\begin{tikzpicture}[
    font=\zihao{-5},
    base/.style={draw=black, line width=0.6pt, align=center, inner sep=4pt, minimum height=1.15cm},
    proc/.style={base, minimum width=#1},
    proc/.default=5.2cm,
    data/.style={base, fill=black!10, minimum width=#1},
    data/.default=5.2cm,
    group/.style={draw=black!60, dashed, line width=0.5pt, inner sep=7pt, minimum width=12.6cm},
    glabel/.style={rotate=90, anchor=south, inner sep=3pt},
    arr/.style={-{Stealth[length=2.2mm, width=1.6mm]}, line width=0.6pt},
    link/.style={line width=0.6pt},
]

% ---------------- 知识与逻辑准备 ----------------
\node[proc] (m1) at (-2.95, 0) {模块一：专家种子构建\\七类业务场景约 600 条种子问答};
\node[data] (corpus) at (2.95, 0) {预处理后的七类电力语料\\（第 2 章）};
\node[proc] (m2) at (-2.95, -1.75) {模块二：设计逻辑提取\\逻辑模板 $\ell = (a, \Xi, \sigma, \phi)$};
\node[proc] (m3) at (2.95, -1.75) {模块三：知识图谱构建\\约 4.2 万条三元组（不含评测文档）};

% ---------------- 生成与演化 ----------------
\node[proc=7.6cm] (m4) at (0, -3.95) {模块四：图遍历指令生成\\沿知识需求链采样图谱路径，按逻辑模板生成问答};
\node[proc=7.6cm] (m5) at (0, -5.6) {模块五：指令演化\\施加五类演化操作提升复杂度，并同步重写答案};

% ---------------- 质量筛选 ----------------
\node[proc=7.6cm] (m6) at (0, -7.8) {模块六：自动化质量筛选\\四维加权评分（$Q \ge 0.75$），一票否决与重复簇去重};
\node[data=7.6cm] (out) at (0, -9.6) {SFT 指令数据集（约 6.8 万条）\\覆盖四类核心任务与三级难度};

% ---------------- 连线 ----------------
\draw[arr] (m1) -- (m2);
\draw[arr] (corpus) -- (m3);
\coordinate (merge) at (0, -2.95);
\draw[link] (m2.south) |- (merge);
\draw[link] (m3.south) |- (merge);
\draw[arr] (merge) -- (m4.north);
\draw[arr] (m4) -- (m5);
\draw[arr] (m5) -- (m6);
\draw[arr] (m6) -- (out);

\draw[arr, dashed] (m6.east) -- ++(0.7, 0) coordinate (fb) |- (m5.east);
\node[font=\zihao{6}, anchor=west, inner sep=2pt] at ($(fb)!0.5!(fb |- m5.east)$) {失败样本反馈};

% ---------------- 分组框 ----------------
\begin{scope}[on background layer]
    \node[group, fit=(m1)(corpus)(m2)(m3)] (g1) {};
    \node[group, fit=(m4)(m5)] (g2) {};
    \node[group, fit=(m6)] (g3) {};
\end{scope}
\node[glabel] at (g1.west) {知识与逻辑准备};
\node[glabel] at (g2.west) {生成与演化};
\node[glabel] at (g3.west) {质量筛选};

\end{tikzpicture}
\end{document}
